\section{Vstavané predikáty}

Klasický Datalog povoľuje prácu iba s predikátmi definovanými v programe, či už pomocou faktov alebo
pravidiel. V prípade, že chceme pracovať s veľkou databázou s niekoľkými tabuľkami, museli by sme
pre každú tabuľku vytvoriť predikát a naplniť ho faktami v každom dotaze. To by bolo veľmi
neefektívne a nepraktické. Preto zavedieme tzv. \textbf{vstavané predikáty}. Tie budú práve
reprezentovať veľké tabuľky alebo predikáty, ktoré by bolo nepraktické definovať a implementovať v
Datalogu. Implementované budú v imperatívnom programovacom jazyku, v našom prípade v Jave, čiže môžu
pracovať aj so zložitejšími vstupmi, napríklad s inými predikátmi, termami a zoznamami. Vieme ich
použiť v pravidlách rovnako ako bežné predikáty, s jedným obmedzením. Nemôžeme ich predefinovať ani
dodefinovať, čiže sa nemôžu vyskytovať v dôsledku implikácie. Vstavané predikáty budeme označovať
pomocou \texttt{\$} pred menom predikátu.

\begin{definition}
    \label{parameter}
    Pod pojmom \textbf{parameter} rozumieme:
    \begin{enumerate}
        \item ľubovoľný term
        \item ľubovoľný predikát alebo vstavaný predikát
        \item ľubovoľný zoznam parametrov, ktorý zapisujeme ako $[P_1, P_2, \dots, P_n]$,\\ kde
        $P_1, P_2, \dots, P_n$ sú parametre
    \end{enumerate}
\end{definition}

\begin{definition}
    \textbf{Atomické formuly} sú
    \begin{enumerate}
        \item $p(t_1, t_2, \dots, t_n)$, kde $p$ je $n$-árny predikát a $t_1, t_2, \dots, t_n$ sú
        termy
        \item $\$p(P_1, P_2, \dots, P_n)$, kde $\$p$ je $n$-árny vstavaný predikát a $P_1, P_2,
        \dots, P_n$ sú parametre.
    \end{enumerate}
\end{definition}
